\section*{Hoofdstuk \ref{ch:g8:protecting-the-nervous-system} --- Het zenuwstelsel beschermen}

\begin{solution}{exo:g8:protecting-the-nervous-system:1}
Ze zet de synapsen vast die overdag versterkt zijn --- het archiveren van
het geleerde --- en draait het chemische huishouden van de \omterm{def:g5:brain-in-command:system}{hersenen}. De
puberende \omterm{def:g5:brain-in-command:system}{hersenen} hebben de uren het hardst nodig: midden in de
herbedrading, en met een klok die naar achteren schuift tegenover een
school die vroeg begint.
\end{solution}

\begin{solution}{exo:g8:protecting-the-nervous-system:2}
Omdat de zintuigontvangers van het \omterm{def:g1:the-five-senses:sense}{oor} per dosis afsterven --- volume maal
blootstelling --- en nooit terug groeien. Uitgegeven: het kapitaal aan
ontvangers. Waarschuwing: het fluiten na een luide avond.
\end{solution}

\begin{solution}{exo:g8:protecting-the-nervous-system:3}
Bij de synapsen --- door een boodschapper na te bootsen, er een te
blokkeren, of afgifte af te dwingen.
\end{solution}

\begin{solution}{exo:g8:protecting-the-nervous-system:4}
Ze dempt de koppelingen in het hele stelsel: het oordeel (het afwegen), de
\omterm{def:g5:brain-in-command:reaction}{reactietijd} (de klok van de lus), het evenwicht en het geheugen --- bij elk
station trager en botter.
\end{solution}

\begin{solution}{exo:g8:protecting-the-nervous-system:5}
Herhaalde doses laten de beloningskoppelingen zich aanpassen --- receptoren
teruggesnoeid, drempels verhoogd --- tot de stof nodig is om je normaal te
voelen. Haar afwezigheid is dan een protest van het stelsel: de uitgang
loopt bergop tegen herbedrade synapsen in.
\end{solution}

\begin{solution}{exo:g8:protecting-the-nervous-system:6}
Omdat ze dezelfde beloningskoppelingen traint, zonder molecuul:
onvoorspelbare kleine prijzen zijn het schema dat bedrading het hardst
versterkt, en de rekening landt op dezelfde \omterm{prop:g1:healthy-habits:needs}{slaap} en aandacht die het
hoofdstuk beschermt.
\end{solution}

\begin{solution}{exo:g8:protecting-the-nervous-system:7}
Leren versterkt koppelingen; \omterm{prop:g1:healthy-habits:needs}{slaap} zet ze vast. Wie van middernacht tot
twee doorging, kocht ruwe oversteken en verkocht het archiveren; wie om
tien uur ging slapen, hield allebei. Stof die niet vastgezet is, is
maar half geleerd.
\end{solution}

\begin{solution}{exo:g8:protecting-the-nervous-system:8}
De puberende \omterm{def:g5:brain-in-command:system}{hersenen} zijn in aanbouw --- verstoorde instructies klinken
door de hele bouw heen (de herbedrading uit
\cref{rem:g8:puberty:moods}) --- en hun koppelingen zijn op hun
instelbaarst, dus zet de beloningsherbedrading van \omterm{prop:g5:healthy-choices:substances}{verslaving} zich juist
daar het snelst vast.
\end{solution}

\begin{solution}{exo:g8:protecting-the-nervous-system:9}
De meters die een auto aflegt vóór de rem zijn \omterm{def:g5:brain-in-command:reaction}{reactietijd} in afstand
uitgedrukt --- en alcohol verlengt de lus bij haar gedempte koppelingen:
melden, afwegen en bevelen worden allemaal trager. De grenzen in de wet
zijn mechanisch gezien een plafond op extra meters.
\end{solution}

\begin{solution}{exo:g8:protecting-the-nervous-system:10}
\omterm{prop:g1:healthy-habits:needs}{Slaap} de uren: het vastzetten loopt 's nachts. Begroot het lawaai:
ontvangers zijn kapitaal dat niet hernieuwbaar is. Geen
koppelingsmisleiders: de stoffen herbedraden instelbare koppelingen
richting de val --- en \omterm{def:g5:brain-in-command:system}{hersenen} die nog in aanbouw zijn herbedraden het
snelst. (Helmen: bedrading geneest niet zoals \omterm{def:g3:bones-and-muscles:skeleton}{bot}; gebruik de machine:
koppelingen worden sterker van verkeer.)
\end{solution}

\begin{solution}{exo:g8:protecting-the-nervous-system:11}
Hulp zoeken wordt onderhoud in plaats van een bekentenis: een
bedradingsprobleem heeft behandelingen, adressen en geen schaamte --- dus
zoeken wie vastzitten eerder hulp. Hulp aanbieden wordt praktisch: het
probleem van een vriend is iets om samen mee naar de deur van de
verpleegkundige te lopen, en geen geheim om te bewaren of een oordeel om
te vellen.
\end{solution}

\begin{solution}{exo:g8:protecting-the-nervous-system:12}
De spleten die zich instellen zijn de reden dat de machine kan leren:
gebruik versterkt oversteken, en vaardigheid, geheugen en elke getrainde
lus zijn vastgezette koppelingen --- de maand van de jongleur. Diezelfde
instelbaarheid is de kwetsbaarheid: stoffen en prijzenschema's die de
beloningskoppelingen doseren, stellen die ook in --- richting het nodig
hebben van de dosis --- en \omterm{prop:g5:healthy-choices:substances}{verslaving} is het mechanisme van het leren, met
de verkeerde les gevoed. Eén eigenschap, twee gezichten: bescherm de
instelbaarheid en ze stapelt zich op tot vaardigheden; geef haar aan
koppelingsmisleiders en ze stapelt zich op tot vallen.
\end{solution}

\begin{solution}{pb:g8:protecting-the-nervous-system:1}
\textbf{1.} Voorspelling: de uitgeruste vrijwilliger vangt centimeters
eerder, en gelijkmatiger over de pogingen heen. Mechanisme: de \omterm{prop:g1:healthy-habits:needs}{slaap} heeft
de koppelingen van die dag vastgezet en de chemische slijtage opgeruimd ---
de lus met een korte nacht weegt bij elke \omterm{def:g8:nervous-communication:synapse}{synaps} traag af.

\textbf{2.} ``De meter laat het dosistempo zien; de doppen snijden het weg.
De zintuigontvangers van het \omterm{def:g1:the-five-senses:sense}{oor} sterven af naar opgetelde dosis en groeien
nooit terug --- gehoor is kapitaal, dat het hardst uitgegeven wordt op
concerten en bij oordopjes op vol volume, en verdedigd wordt met zachter,
korter en met doppen.''

\textbf{3.} Nabootsen: nicotine --- die zich bij de beloningskoppelingen
voordoet als een natuurlijke boodschapper. Blokkeren: het verlammende gif
--- de boodschapper bij de spierkoppeling geweigerd. Forceren: alcohol
past in het schema als de demper van het hele stelsel (zijn pijl: het
gesprek van de koppelingen van begin tot eind dempen).

\textbf{4.} Het doelwit: dezelfde beloningskoppelingen, getraind door
hetzelfde schema van onvoorspelbare prijzen --- de meldingendemonstratie is
een doseerdemonstratie met licht en geluid in plaats van moleculen.

\textbf{5.} Op je veertiende worden de koppelingen nog gebouwd en zijn ze
op hun instelbaarst: dezelfde dosis herbedraadt dieper en zet zich sneller
vast dan in afgebouwde \omterm{def:g5:brain-in-command:system}{hersenen} --- ``het is goed met ons gekomen'' werd op een
lagere moeilijkheidsgraad gespeeld.

\textbf{6.} De prijs van de eerste sigaret is die van alle volgende: elke
dosis begint de \omterm{def:g4:adapted-to-their-world:adaptation}{aanpassing} van de beloningskoppelingen, en de ingang is
precies één beslissing breed --- de uitgang jaren herbedraad protest. Wat
aangeboden wordt is de deuropening van de val, geen sigaret.

\textbf{7.} De televisie plande haar beloningen; de tijdlijn deelt ze
onvoorspelbaar uit, item voor item --- en onvoorspelbare kleine prijzen
zijn het schema dat koppelingen het hardst traint, en daarom verzet dat
checken zich tegen stoppen zoals programma's dat nooit deden.

\textbf{8.} Dat alleen stoppen vechten is tegen herbedrade koppelingen, en
dat niemand dat hoeft: verpleegkundigen en artsen behandelen het
bedradingsprobleem vertrouwelijk --- ook bij minderjarigen --- en opnieuw
aanpassen werkt. Eerste stap: vandaag één volwassene die je vertrouwt,
want zwijgen is de enige zet die verliest.

\textbf{9.} ``Een maand jongleren versterkt koppelingen tot drie ballen
zichzelf vliegen --- dat is leren, en daar zijn je synapsen voor.
\omterm{prop:g5:healthy-choices:substances}{Verslaving} is precies hetzelfde mechanisme: doses versterken de verkeerde
lus tot die zichzelf draait. De \omterm{def:g5:brain-in-command:system}{hersenen} kunnen het ene niet bieden zonder
het andere te riskeren --- de instelbaarheid ís het geschenk. Dus is de
vraag nooit of jouw koppelingen getraind zullen worden, maar alleen
waardoor.''

\textbf{10.} Omdat het moment van het aanbod het slechtste moment is om te
gaan afwegen --- de druk belast juist de koppelingen die het beslissen
doen. Een zin die je rustig van tevoren bedenkt, maakt van de weigering een
reflex: voorafgewogen routering, klaar voordat ze nodig is.

\textbf{11.} Een verdedigbare volgorde: \omterm{prop:g1:healthy-habits:needs}{slaap} eerst (die onderhoudt de
winst van al het andere), dan geen koppelingsmisleiders (de onomkeerbare
val), dan het lawaai (onvervangbaar kapitaal), dan de schermuren, dan de
helmen. De eerste keuze verdedigd: de opbrengst van elk ander punt wordt
's nachts vastgezet --- \omterm{prop:g1:healthy-habits:needs}{slaap} is het onderhoud waar de rest van afhangt.

\textbf{12.} ``We beschermen de instelbare koppelingen die alles zullen
hebben geleerd wat jij ooit weet --- en veertien zit midden in hun
snelste, best trainbare en makkelijkst te vangen jaren.''
\end{solution}
